\documentclass[11pt]{article}

\usepackage[margin=1in]{geometry}
\usepackage{amsmath,amssymb,amsthm}
\usepackage{enumitem}
\usepackage[hidelinks]{hyperref}

\newtheorem{theorem}{Theorem}[section]
\newtheorem{lemma}[theorem]{Lemma}
\newtheorem{proposition}[theorem]{Proposition}
\newtheorem{corollary}[theorem]{Corollary}
\newtheorem{conjecture}[theorem]{Conjecture}
\newtheorem{definition}[theorem]{Definition}

\theoremstyle{remark}
\newtheorem{remark}[theorem]{Remark}

\newcommand{\F}{\mathbb F}
\newcommand{\Z}{\mathbb Z}
\newcommand{\one}{\mathbf 1}
\newcommand{\supp}{\operatorname{supp}}

\title{Dyadic tail obstructions for regular induced subgraphs}
\author{}
\date{}

\begin{document}
\maketitle

\begin{abstract}
Let \(F(n)\) be the largest integer \(k\) such that every graph on \(n\)
vertices contains a regular induced subgraph on at least \(k\) vertices.  The
conjecture of Erdos, Fajtlowicz, and Staton asks whether
\(F(n)/\log n\to\infty\).  This manuscript gives a self-contained account of a
dyadic reduction of the conjecture to one finite first-return statement.  The
reduction is unconditional; the final superlogarithmic conclusion is stated
conditionally on a q-marker provenance/support-decrease theorem.  We prove the
terminal congruence criterion, the complement-pair carry identity, the
first-return marker reduction, the local marker exits, the product-firewall
normal form, and the polynomial-threshold consequence of the missing
provenance theorem.  The point is to make the remaining mathematical gap
explicit: an ambient splitter of a fully skew q-marker must be routed to an
ordered first-return boundary row or to a local regularizing exit.
\end{abstract}

\tableofcontents

\section{Introduction}

All graphs in this paper are finite, simple, and undirected.  If
\(S\subseteq V(G)\), we write \(G[S]\) for the subgraph induced by \(S\) and
\(\deg_S(v)\) for the degree of \(v\) inside \(G[S]\).

Define
\[
  f(G)=\max\{|S|:G[S]\text{ is regular}\},\qquad
  F(n)=\min_{|V(G)|=n} f(G).
\]
The classical Ramsey bound gives \(F(n)\ge c\log n\).  Erdos,
Fajtlowicz, and Staton asked whether this lower bound is always
superlogarithmic:
\[
  \frac{F(n)}{\log n}\longrightarrow\infty .
\]
This remains the global problem.  The present manuscript does not claim an
unconditional proof of the conjecture.  Instead, it isolates the exact finite
statement to which the dyadic approach reduces it and proves the surrounding
machinery in full.

The strategy is modular.  A set \(U\) of size \(q\) is regular as soon as all
degrees in \(G[U]\) are congruent modulo \(q\), because the degrees already lie
in the interval \([0,q-1]\).  Thus the terminal problem is to force congruence
of induced degrees modulo \(q\).  In a dyadic descent, \(q=2^j\), one reaches a
terminal host \(U\) with a dropped tail \(R\).  The only remaining datum is
\[
  \rho_R(u)=|N(u)\cap R|,\qquad u\in U.
\]
If \(\rho_R\) is constant modulo \(q\), then \(G[U]\) is regular.  The possible
failure at the next dyadic bit is measured by a quotient class
\(\tau_m(R,U)\in \F_2^U/\langle\one_U\rangle\).  We identify this class with an
aggregate complement-orbit carry \(\beta_m\).  A nonzero \(\beta_m\) gives a
first-return marker.

The remaining local problem is the q-marker endpoint.  The low-set update
proves that every such marker has size \(0\) modulo \(q\), so any nonempty
submarker of size \(<q\) is impossible.  Many local alternatives are eliminated
directly: fixed-trace exits, clean marked exits, frozen same-trace markers,
one-splitter markers, and the static selections in the fully skew split
quotient.  The survivor is a product firewall: provenance rows cut only the
residual-pair/compensator coordinate, while ambient primeness cuts a protected
packet only through high-error breakers.  The missing theorem is precisely the
claim that such ambient breakers can be promoted to ordered first-return
boundary rows, or else cause local exits.  Once that theorem is supplied, a
sub-\(q\) marker appears and the endpoint collapses.

\section{Background and threshold form}

Let
\[
  T(k)=\min\{n:\text{ every graph on }n\text{ vertices contains a regular
  induced subgraph on }k\text{ vertices}\}.
\]
Then \(k\le F(n)\) if and only if \(T(k)\le n\).

\begin{lemma}[threshold reduction]\label{lem:threshold}
Assume that there is a constant \(A\) such that
\[
  T(2^r)\le 2^{A(r+1)^2}
\]
for all sufficiently large \(r\).  Then \(F(n)/\log n\to\infty\).
\end{lemma}

\begin{proof}
It is enough to prove that, for every fixed \(M\), all sufficiently large powers
\(2^k\) have \(F(2^k)\ge Mk\).  Choose \(r=r(k)\) maximal with
\[
  A(r+1)^2\le k.
\]
Then \(r\to\infty\), and hence \(2^r\ge Mk\) for all sufficiently large \(k\).
The assumed threshold bound gives
\[
  T(2^r)\le 2^{A(r+1)^2}\le 2^k,
\]
so \(F(2^k)\ge 2^r\ge Mk\).

For general \(n\), take \(k\) with \(2^k\le n<2^{k+1}\).  Monotonicity gives
\(F(n)\ge F(2^k)\), while \(\log n\le (k+1)\log 2\).  Since \(M\) was
arbitrary, the ratio \(F(n)/\log n\) tends to infinity.
\end{proof}

\section{Modular terminal sets}

\begin{definition}
Let \(q\ge 1\).  A vertex set \(U\) is \(q\)-modular if
\[
  \deg_U(u)\equiv\deg_U(v)\pmod q
\]
for all \(u,v\in U\).
\end{definition}

\begin{lemma}[terminal congruence collapse]\label{lem:terminal-collapse}
If \(|U|\le q\) and \(U\) is \(q\)-modular, then \(G[U]\) is regular.
\end{lemma}

\begin{proof}
Every degree in \(G[U]\) lies in \(\{0,\ldots,|U|-1\}\), an interval of length
strictly less than \(q\).  Two integers in such an interval that are congruent
modulo \(q\) are equal.
\end{proof}

\begin{definition}
Let \(U\) and \(R\) be disjoint vertex sets.  We call \((U,R)\) a
fixed-modulus terminal host at modulus \(q\) if \(|U|=q\) and the degrees in
\(G[U\cup R]\) are constant modulo \(q\) on \(U\).  Write
\[
  \rho_R(u)=|N(u)\cap R|,\qquad u\in U.
\]
\end{definition}

\begin{lemma}[tail criterion]\label{lem:tail-criterion}
Let \((U,R)\) be a fixed-modulus terminal host at modulus \(q\).  If
\(\rho_R\) is constant modulo \(q\) on \(U\), then \(G[U]\) is regular.
\end{lemma}

\begin{proof}
For \(u,v\in U\),
\[
  \deg_{U\cup R}(u)-\deg_{U\cup R}(v)
  =\deg_U(u)-\deg_U(v)+\rho_R(u)-\rho_R(v).
\]
The first difference is \(0\) modulo \(q\) by the host condition, and the tail
difference is \(0\) modulo \(q\) by hypothesis.  Thus
\(\deg_U(u)\equiv\deg_U(v)\pmod q\).  Lemma~\ref{lem:terminal-collapse}
applies because \(|U|=q\).
\end{proof}

\section{Dyadic tail classes}

Throughout this section \(q=2^j\) and \((U,R)\) is a terminal host.  Suppose the
first \(m\) dyadic bits of the tail have been frozen:
\[
  \rho_R=K_m\one_U+2^m h_m,\qquad m<j.
\]
Define the terminal-tail class
\[
  \tau_m(R,U)=[h_m\bmod 2]\in \F_2^U/\langle\one_U\rangle.
\]
The next bit is frozen exactly when \(\tau_m(R,U)=0\).

For each \(A\subseteq U\), let
\[
  n_A=|\{x\in R:N(x)\cap U=A\}|.
\]
The complement pair of \(A\) is \([A]=\{A,U\setminus A\}\).  Define
\[
  \epsilon_{[A]}
  =\left\lfloor\frac{n_A+n_{U\setminus A}}{2^m}\right\rfloor\pmod 2
\]
and
\[
  \beta_m=\sum_{[A]}\epsilon_{[A]}\,[1_A]
  \in \F_2^U/\langle\one_U\rangle.
\]

\begin{proposition}[aggregate carry identity]\label{prop:beta}
For every \(m<j\),
\[
  \tau_m(R,U)=\beta_m.
\]
\end{proposition}

\begin{proof}
For \(u,v\in U\),
\[
  \rho_R(u)-\rho_R(v)
  =\sum_{A\subseteq U} n_A\bigl(1_A(u)-1_A(v)\bigr).
\]
Inside the quotient by constants,
\[
  1_{U\setminus A}=\one_U-1_A\equiv -1_A.
\]
Thus the two members of the complement pair \([A]\) contribute to pairwise
differences through the single coefficient \(n_A+n_{U\setminus A}\).  Dividing
by \(2^m\) and reducing modulo \(2\), the only surviving contribution of the
pair is
\[
  \left\lfloor\frac{n_A+n_{U\setminus A}}{2^m}\right\rfloor
  \bigl(1_A(u)-1_A(v)\bigr)\pmod 2.
\]
Pairwise differences determine a vector modulo constants, so the class of
\(h_m\bmod2\) is exactly the displayed aggregate \(\beta_m\).
\end{proof}

\begin{corollary}[cut extraction]\label{cor:cut}
If \(\beta_m\ne0\), then there is a nonempty proper cut \(S\subset U\), unique
up to complement after merging equal terminal profiles, such that the normalized
next-bit cocycle
\[
  \kappa_m(u,v)=\frac{\rho_R(u)-\rho_R(v)}{2^m}\pmod 2
\]
is \(1\) across \(S\) and \(0\) within each side.
\end{corollary}

\begin{proof}
Choose a representative \(b\in\F_2^U\) of the nonzero quotient class
\(\beta_m\).  Its support is nonempty and proper, after possibly replacing
\(b\) by \(b+\one_U\).  Proposition~\ref{prop:beta} identifies this
representative with the normalized next-bit differences of \(\rho_R\).  Merging
terminal profile classes on which \(b\) is constant makes the cut unique up to
complement.
\end{proof}

\section{First-return markers}

The next definitions encode the local object produced by a nonzero cut.

\begin{definition}
A first-return boundary is fixed-trace if it preserves the trace of the current
cut, clean if it changes only one already marked coordinate, and mixed
otherwise.
\end{definition}

\begin{lemma}[fixed and clean boundaries]\label{lem:fixed-clean}
In a minimal obstruction to vanishing of \(\beta_m\), the first boundary seeing
the cut of Corollary~\ref{cor:cut} is mixed.
\end{lemma}

\begin{proof}
If the first boundary is fixed-trace, then all vertices in the corresponding
trace class have the same outside contribution.  If that class contains an
induced \(P_3\), the middle vertex and an endpoint have equal outside trace and
internal degrees differing by one.  The first-return exchange supplies the
opposite unit discrepancy, giving a regular \(q\)-set.  If there is no induced
\(P_3\), each connected component is a clique: a shortest path between two
nonadjacent vertices would otherwise contain an induced \(P_3\).  Equal outside
trace and the modular equations force the relevant clique sizes to have a common
residue modulo \(q\); selecting equal portions from enough clique components
gives a regular \(q\)-set.

If the boundary is clean, the two repairs differ by opposite unit vectors on a
single marked coordinate.  Their sum is zero in the quotient by constants, so
one repair preserves the obstruction on a smaller profile suborbit, contrary to
minimality.
\end{proof}

\begin{definition}[q-marker]
A q-marker is a tuple \((d,e,x,y;R)\), with \(R\ne\varnothing\), such that:
\begin{enumerate}[label=(\roman*)]
\item \(d,e\) have the same trace outside \(R\);
\item \(x,y\) have the same trace outside \(R\);
\item the two traces differ exactly on \(R\);
\item the marked quartet is one of the two complementary two-class patterns
      (miss/miss or add/add);
\item \(|R|\equiv0\pmod q\).
\end{enumerate}
\end{definition}

\begin{lemma}[mixed boundary to marker]\label{lem:marker}
A minimal mixed first-return boundary produces a q-marker.
\end{lemma}

\begin{proof}
Let \(d,e\) be the deleted pair and \(x,y\) the inserted pair in the first failed
exchange.  First-return minimality makes the deleted pair one trace class and
the inserted pair another trace class outside the slack rows spent by the
exchange.  Let \(R\) be exactly the set of slack rows on which the two traces
differ.  All modular degree equations are preserved except for the contribution
of \(R\).  Since the failed exchange occurs inside a fixed-modulus host, that
contribution is \(0\) modulo \(q\), so \(|R|\equiv0\pmod q\).  The two signs of
the exchange give the miss/miss and add/add quartet patterns.
\end{proof}

\section{Local q-marker exits}

\begin{lemma}[no-split module exit]\label{lem:module}
If no outside row splits the marker set \(R\), then \(R\) is a module.  Hence no
minimal prime-shell counterexample has an unsplit q-marker.
\end{lemma}

\begin{proof}
If no outside row splits \(R\), every outside vertex is complete or
anticomplete to \(R\), which is the definition of a module.  In a minimal
counterexample, a nontrivial module can be contracted or replaced by an
equally-sized same-internal-degree subset without changing any outside
adjacency.  Thus such a module is impossible.
\end{proof}

\begin{lemma}[frozen marker closure]\label{lem:frozen}
Suppose every splitter is constant on each same-trace component of a q-marker.
Then the graph contains a regular induced \(q\)-set.
\end{lemma}

\begin{proof}
Within a same-trace component all outside contributions are equal.  If one such
component contains an induced \(P_3\), the middle vertex and an endpoint differ
by one internal neighbor and have the same outside trace.  The marked quartet
supplies the opposite unit discrepancy, giving a regular \(q\)-set by exchange.

If no same-trace component contains an induced \(P_3\), then every component is
a clique.  Since \(|R|\equiv0\pmod q\), the clique sizes in a frozen trace fiber
are congruent modulo \(q\).  If a clique has size at least \(q\), it contains
\(K_q\).  Otherwise all relevant clique sizes are some \(s<q\).  If
\(g=\gcd(s,q)\), then the number of cliques is divisible by \(q/g\), and taking
\(g\) vertices from each of \(q/g\) cliques gives an induced \(q\)-vertex
\((g-1)\)-regular graph.
\end{proof}

\begin{lemma}[one-splitter closure]\label{lem:one-splitter}
If a minimal q-marker has a splitter but no second independent splitter, then
there is a regular induced \(q\)-set or a smaller q-marker.
\end{lemma}

\begin{proof}
Let \(z\) be the unique splitter up to complement of its cut on \(R\).  The
split gives a one-excess profile across the two sides.  If a neighbor of \(z\)
is isolated inside the marker side, deleting that vertex leaves a balanced
regular \(q\)-set.  Otherwise every neighbor of \(z\) is supported by an
internal marker edge, so an alternating length-two path appears.  Together with
the marked quartet this is the same false-clone \(P_3\) exchange as in
Lemma~\ref{lem:frozen}, and it produces a smaller marker or a local regularizing
exit.
\end{proof}

\section{The fully skew split quotient}

After the local exits above, a surviving q-marker must be fully skew.  The
finite quotient arithmetic gives the following endpoint.

\begin{proposition}[split quotient]\label{prop:split-quotient}
In the fully skew endpoint, after removing local exits and proper submarkers,
the active marker graph has the form
\[
  R=A\sqcup U,\qquad A=K_{q-2},\qquad U=\{u_0,u_1\},
\]
and \(G[U]\) is either \(K_2\) or \(2K_1\).  If
\[
  k(H)=
  \begin{cases}
    (q-4)/2,&H=K_2,\\
    q/2-1,&H=2K_1,
  \end{cases}
\]
then every regular selection using \(U\) requires one compensator component
containing a \(k(H)\)-clique.  A selection avoiding \(U\) is exactly a
proper-divisor equal-clique bypass.
\end{proposition}

\begin{proof}
If a residual vertex has mixed adjacency to \(A\), two adjacent vertices of
\(A\) together with it give a same-trace \(P_3\), already handled by
Lemma~\ref{lem:frozen}.  If a residual vertex is complete to \(A\), then adding
one inserted root gives a \(K_q\) in the relevant orientation.  Hence the
residual pair is anticomplete to \(A\), and \(G[U]\) is \(K_2\) or \(2K_1\).

A regular selection avoiding \(U\) sees only clique packets, so equal degree
requires equal-size choices from equal-trace cliques; these are exactly the
proper-divisor equal-clique selections.  If the selection uses \(U\), the
residual pair has equal marked trace, and its degree deficit must be supplied by
a one-sided compensator.  Direct degree equality gives deficit \((q-4)/2\) in
the \(K_2\) case and \(q/2-1\) in the \(2K_1\) case.  Vertices selected from
different smaller compensator components are nonadjacent while having the same
outside trace, so spreading the compensator over several smaller components
cannot equalize degrees.  Thus a single \(k(H)\)-clique is necessary.
\end{proof}

\section{The missing provenance theorem}

The endpoint left by Proposition~\ref{prop:split-quotient} is a product
firewall.  Provenance rows cut only the residual pair or its compensator, while
ambient primeness may split \(A\) or a compensator component only by a
high-error breaker.  The following theorem is the missing finite input.

\begin{conjecture}[q-marker provenance/support-decrease]\label{conj:qmarker}
In a fully skew first-return q-marker endpoint, every ambient splitter of a
protected packet has one of the following outcomes:
\begin{enumerate}[label=(\roman*)]
\item it is promoted to an ordered first-return boundary row;
\item it gives a local regularizing exit;
\item its first failed row carries a complete smaller first-return q-marker.
\end{enumerate}
Moreover, in case (iii) the smaller marker is marker-complete: it is the whole
shared-slack set of a new marked two-class quartet, not merely a failed prefix.
\end{conjecture}

\begin{remark}
This is the point at which the current proof route is conditional.  Primeness
alone gives an arbitrary splitter of the marker, but not necessarily a splitter
with first-return provenance.  Static degree balancing can also create splitters
without square-provenance.  Thus Conjecture~\ref{conj:qmarker} is not a
bookkeeping statement; it is the remaining ordered first-return theorem.
\end{remark}

\begin{lemma}[ordered boundary completeness]\label{lem:boundary-complete}
Assume an ambient high-error packet breaker has already been promoted to an
ordered boundary row \(z\).  If \(z\) splits a protected packet \(P\), then the
first packet-internal failed set is a whole side of the split,
\[
  P\cap N(z)\quad\text{or}\quad P\setminus N(z).
\]
In particular, if \(P\) is proper, this failed set has size \(<q\).
\end{lemma}

\begin{proof}
All vertices of \(P\) have the same trace to the marked quartet and to every
packet other than \(P\).  The ordered boundary test can distinguish vertices of
\(P\) only through adjacency to \(z\).  Hence the failure indicator is constant
on \(P\cap N(z)\) and on \(P\setminus N(z)\).  First failure cannot select a
proper nonempty subset of one side, because an omitted and an included vertex on
the same side have identical boundary data.  Therefore the failed set is a whole
side.  The protected packets in the fully skew endpoint are \(A\) of size
\(q-2\) and compensator components of size \(<k(H)<q\), so every proper side has
size \(<q\).
\end{proof}

\begin{proposition}[conditional q-marker exclusion]\label{prop:conditional-marker}
Assume Conjecture~\ref{conj:qmarker}.  Then no minimal fully skew q-marker
survives after the local exits of Lemmas~\ref{lem:module},
\ref{lem:frozen}, and \ref{lem:one-splitter}.
\end{proposition}

\begin{proof}
Let a fully skew survivor be chosen minimally.  By
Proposition~\ref{prop:split-quotient}, all static regular selections have been
removed and the survivor has product-firewall form.  Ambient primeness supplies
a high-error breaker of a protected packet.  Apply
Conjecture~\ref{conj:qmarker}.  A local exit contradicts survival.  A smaller
complete q-marker contradicts minimality.  If the breaker is promoted to an
ordered boundary row, Lemma~\ref{lem:boundary-complete} says that the first
packet-internal failed set is a whole side of a proper protected packet, hence
nonempty of size \(<q\).  The low-set update makes that side a q-marker, so its
size is \(0\) modulo \(q\), impossible for a nonempty set of size \(<q\).
\end{proof}

\section{Conditional terminal theorem}

\begin{theorem}[conditional dyadic terminal theorem]\label{thm:conditional-terminal}
Assume Conjecture~\ref{conj:qmarker}.  Let \(q=2^j>1\) and let \((U,R)\) be a
minimal terminal host at modulus \(q\).  Then \(\rho_R\) is constant modulo
\(q\), and \(G\) contains a regular induced \(q\)-set.
\end{theorem}

\begin{proof}
Suppose \(\rho_R\) is not constant modulo \(q\).  Let \(m\) be the first dyadic
bit at which constancy fails.  Then \(\tau_m(R,U)\ne0\), so
\(\beta_m\ne0\) by Proposition~\ref{prop:beta}.  Corollary~\ref{cor:cut}
produces a nonconstant cut.  Lemma~\ref{lem:fixed-clean} shows that the first
boundary seeing this cut is mixed, and Lemma~\ref{lem:marker} turns the mixed
boundary into a q-marker.  The local exits and
Proposition~\ref{prop:conditional-marker} exclude the marker, contradiction.
Thus every dyadic bit is frozen and \(\rho_R\) is constant modulo \(q\).  The
regular \(q\)-set follows from Lemma~\ref{lem:tail-criterion}.
\end{proof}

\section{Polynomial dyadic lifting}

\begin{lemma}[parity base]\label{lem:parity}
Every graph on \(n\) vertices has a modulus-\(2\) cascade witness of order at
least \(n/2\).
\end{lemma}

\begin{proof}
Pair vertices greedily, leaving at most one unpaired vertex.  On each pair, the
induced degrees are equal modulo \(2\): adjacent pairs have degrees \(1,1\), and
nonadjacent pairs have degrees \(0,0\).  The collection of pairs is the base
modulus-\(2\) witness.
\end{proof}

\begin{lemma}[conditional dyadic lift]\label{lem:lift}
Assume Conjecture~\ref{conj:qmarker}.  There is a constant \(C\) such that, for
every dyadic \(q=2^j\), a fixed-modulus cascade witness of order \(q^C m\) at
modulus \(q\) contains a fixed-modulus cascade witness of order \(m\) at modulus
\(2q\).
\end{lemma}

\begin{proof}
Only finitely many residue functions are used in one terminal descent step:
ambient degree modulo \(q\), dropped-tail degree modulo \(q\), and the residues
against the separated control blocks.  Let \(C\) dominate the number of such
functions in the finite exchange templates.  By the pigeonhole principle, a
witness of order \(q^C m\) contains a subwitness of order at least \(m\) on
which all these residues are fixed.  On each terminal bucket, Theorem
\ref{thm:conditional-terminal} freezes the dropped-tail residue modulo \(q\),
which is exactly the missing next dyadic bit.  Thus the selected subwitness is
a modulus-\(2q\) witness.
\end{proof}

\begin{theorem}[conditional threshold]\label{thm:conditional-threshold}
Assume Conjecture~\ref{conj:qmarker}.  There is a constant \(A\) such that
\[
  T(2^r)\le 2^{A(r+1)^2}
\]
for all \(r\).
\end{theorem}

\begin{proof}
Starting from Lemma~\ref{lem:parity}, apply Lemma~\ref{lem:lift} successively
from \(2\) to \(2^r\).  The cost of the step \(2^i\to2^{i+1}\) is at most
\((2^i)^C\).  Hence the total lifting cost is
\[
  \prod_{i=1}^{r-1}(2^i)^C=2^{C r(r-1)/2}.
\]
Multiplying by the terminal target size \(2^r\) and increasing the constant
gives the asserted bound.
\end{proof}

\begin{corollary}[conditional superlogarithmic conclusion]
Assume Conjecture~\ref{conj:qmarker}.  Then
\[
  \frac{F(n)}{\log n}\to\infty .
\]
\end{corollary}

\begin{proof}
Combine Theorem~\ref{thm:conditional-threshold} with
Lemma~\ref{lem:threshold}.
\end{proof}

\section{What remains}

The preceding sections show that the dyadic route no longer lacks asymptotic
bookkeeping.  The unresolved point is finite and local:
Conjecture~\ref{conj:qmarker}.  Equivalently, one must prove
ambient-to-provenance routing for a fully skew q-marker.  The known local
closures show that fixed-trace failures, clean marked failures, no-split
markers, frozen marker components, one-splitter markers, and all static
selections in the \(K_{q-2}\sqcup U\) quotient do not survive.  A remaining
counterexample must keep an ambient high-error packet breaker and a provenance
residual-pair cut in different coordinates.  Proving that first-return order
forces these coordinates into one square-provenance frame is exactly the
remaining theorem.

\begin{thebibliography}{9}

\bibitem{AKS}
N. Alon, M. Krivelevich, and B. Sudakov,
\emph{Large nearly regular induced subgraphs},
SIAM J. Discrete Math. \textbf{22} (2008), no. 4, 1325--1337.

\bibitem{FMRS}
S. Fajtlowicz, C. McColgan, T. Reid, and P. Staton,
\emph{Ramsey numbers for induced regular subgraphs},
Ars Combin. \textbf{39} (1995), 149--154.

\bibitem{KSW}
M. Krivelevich, B. Sudakov, and N. Wormald,
\emph{Regular induced subgraphs of a random graph},
Random Structures Algorithms \textbf{38} (2011), no. 3, 235--250.

\bibitem{Ramsey}
F. P. Ramsey,
\emph{On a problem of formal logic},
Proc. London Math. Soc. \textbf{30} (1930), 264--286.

\end{thebibliography}

\end{document}
